\documentclass[11pt,a4paper]{article}
\usepackage[utf8]{inputenc}
\usepackage[T1]{fontenc}
\usepackage[margin=2.5cm]{geometry}
\usepackage{xcolor}
\usepackage{listings}
\usepackage{longtable}
\usepackage{booktabs}
\usepackage{titlesec}
\usepackage{hyperref}
\usepackage{fancyhdr}

\definecolor{primaryblue}{RGB}{46,82,102}
\definecolor{codebg}{RGB}{242,242,242}

\titleformat{\section}{\Large\bfseries\color{primaryblue}}{\thesection.}{0.5em}{}
\titleformat{\subsection}{\large\bfseries\color{primaryblue}}{\thesubsection}{0.5em}{}

\lstset{
  backgroundcolor=\color{codebg},
  basicstyle=\ttfamily\small,
  breaklines=true,
  frame=none,
  xleftmargin=0.5em,
  xrightmargin=0.5em,
  aboveskip=8pt,
  belowskip=8pt,
  keywordstyle=\color{primaryblue}\bfseries,
  stringstyle=\color{teal}
}

\pagestyle{fancy}
\fancyhf{}
\fancyhead[L]{\small\color{gray}Projet 7 --- CMS pour blog personnel avec MongoDB}
\fancyfoot[C]{\thepage}
\renewcommand{\headrulewidth}{0.4pt}

\hypersetup{colorlinks=true, linkcolor=primaryblue, urlcolor=primaryblue}

\title{\color{primaryblue}\textbf{Modèle de données JSON}}
\author{}
\date{}

\begin{document}
\maketitle
\thispagestyle{fancy}

\noindent Ce document présente la structure JSON de chacune des trois collections MongoDB utilisées par l'application : \texttt{users}, \texttt{articles} et \texttt{comments}. Pour chaque collection, un tableau détaille les champs (nom, type BSON, contraintes, description) suivi d'un exemple de document réel tel qu'il serait stocké en base.

\section{Collection : users}
Représente un compte utilisateur, avec un rôle unique parmi \texttt{ADMIN}, \texttt{AUTHOR} ou \texttt{READER}.

\begin{longtable}{p{2cm}p{2cm}p{2.8cm}p{5cm}}
\toprule
\textbf{Champ} & \textbf{Type BSON} & \textbf{Contrainte} & \textbf{Description} \\
\midrule
\endhead
\texttt{\_id} & ObjectId & Généré automatiquement & Identifiant unique MongoDB du document \\
\texttt{fullName} & String & Requis & Nom complet de l'utilisateur \\
\texttt{email} & String & Requis, unique (index) & Adresse email, utilisée pour l'authentification \\
\texttt{password} & String & Requis & Mot de passe hashé (BCrypt), jamais stocké en clair \\
\texttt{role} & String (enum) & ADMIN \textbar\ AUTHOR \textbar\ READER & Rôle unique déterminant les permissions \\
\texttt{createdAt} & Date (ISODate) & Généré à la création & Date de création du compte \\
\bottomrule
\end{longtable}

\noindent\textbf{Exemple de document}
\begin{lstlisting}
{
  "_id": ObjectId("6a5248391e987ba0b1f5fe41"),
  "fullName": "Test Author",
  "email": "author@test.com",
  "password": "$2a$10$N9qo8uLOickgx2ZMRZoMy.Ig...",
  "role": "AUTHOR",
  "createdAt": ISODate("2026-07-11T13:48:12.000Z")
}
\end{lstlisting}

\section{Collection : articles}
Représente un article de blog. L'auteur et les images sont référencés (\texttt{authorId}, URLs MinIO) ; les commentaires ne sont pas embarqués (voir le rapport sur la flexibilité et les performances).

\begin{longtable}{p{2cm}p{2cm}p{2.8cm}p{5cm}}
\toprule
\textbf{Champ} & \textbf{Type BSON} & \textbf{Contrainte} & \textbf{Description} \\
\midrule
\endhead
\texttt{\_id} & ObjectId & Généré automatiquement & Identifiant unique MongoDB du document \\
\texttt{title} & String & Requis & Titre de l'article \\
\texttt{content} & String & Requis & Contenu textuel de l'article \\
\texttt{authorId} & String & Requis (index) & Référence vers \texttt{\_id} du document \texttt{users}, auteur de l'article \\
\texttt{tags} & Array\textless String\textgreater & Au moins un tag (index multikey) & Liste des mots-clés associés à l'article \\
\texttt{images} & Array\textless String\textgreater & Optionnel & URLs des images stockées sur MinIO \\
\texttt{status} & String (enum) & DRAFT \textbar\ PUBLISHED & État de publication de l'article \\
\texttt{views} & Number (int) & Défaut : 0 & Compteur de vues de l'article \\
\texttt{createdAt} & Date (ISODate) & Généré à la création & Date de création de l'article \\
\texttt{updatedAt} & Date (ISODate) & Mis à jour à chaque modification & Date de dernière modification \\
\bottomrule
\end{longtable}

\noindent\textbf{Exemple de document}
\begin{lstlisting}
{
  "_id": ObjectId("6a524a7b1e987ba0b1f5fe42"),
  "title": "Mon premier article",
  "content": "Contenu de test sur MongoDB et Spring Boot.",
  "authorId": "6a5248391e987ba0b1f5fe41",
  "tags": ["mongodb", "spring", "backend"],
  "images": [
    "http://localhost:9000/blog-images/articles/3f2a-photo.jpg"
  ],
  "status": "PUBLISHED",
  "views": 12,
  "createdAt": ISODate("2026-07-11T13:51:55.952Z"),
  "updatedAt": ISODate("2026-07-11T14:05:10.310Z")
}
\end{lstlisting}

\section{Collection : comments}
Représente un commentaire, rattaché à un article via \texttt{articleId} (référencement, collection séparée). Le nom de l'auteur est dénormalisé (\texttt{authorName} en texte libre) plutôt que référencé vers \texttt{users}, conformément au choix de favoriser la disponibilité en lecture propre au NoSQL.

\begin{longtable}{p{2cm}p{2cm}p{2.8cm}p{5cm}}
\toprule
\textbf{Champ} & \textbf{Type BSON} & \textbf{Contrainte} & \textbf{Description} \\
\midrule
\endhead
\texttt{\_id} & ObjectId & Généré automatiquement & Identifiant unique MongoDB du document \\
\texttt{articleId} & String & Requis (index) & Référence vers \texttt{\_id} du document \texttt{articles} concerné \\
\texttt{authorName} & String & Requis & Nom affiché de l'auteur du commentaire (dénormalisé) \\
\texttt{content} & String & Requis & Contenu du commentaire \\
\texttt{createdAt} & Date (ISODate) & Généré à la création & Date de publication du commentaire \\
\bottomrule
\end{longtable}

\noindent\textbf{Exemple de document}
\begin{lstlisting}
{
  "_id": ObjectId("6a52501e1e987ba0b1f5fe45"),
  "articleId": "6a524a7b1e987ba0b1f5fe42",
  "authorName": "Visiteur",
  "content": "Super article, merci pour le partage !",
  "createdAt": ISODate("2026-07-11T14:12:03.221Z")
}
\end{lstlisting}

\section{Relations entre collections}
Le schéma ci-dessous résume les relations de référencement entre les trois collections~:

\begin{lstlisting}
users (1) -----< articles       (via articles.authorId)
articles (1) ---< comments      (via comments.articleId)
\end{lstlisting}

\noindent Aucune relation n'est modélisée par embarquement dans ce projet : chaque collection reste indépendante et interrogeable individuellement, ce qui correspond aux choix détaillés dans le rapport sur la flexibilité et les performances.

\end{document}
